\documentclass[10pt,twocolumn]{article}

% ----------------------------------------------------------------------------
% INV-15: A Formal Safety Invariant for KV-Cache Reuse in
% Multi-Agent Judge Pipelines
%
% Target venues: arXiv (cs.LG / cs.DC), MLSys 2027,
% NeurIPS 2026 Workshop on Efficient LLM Inference
% ----------------------------------------------------------------------------

\usepackage[a4paper,margin=0.85in,columnsep=0.30in]{geometry}
\usepackage[utf8]{inputenc}
\usepackage[T1]{fontenc}
\usepackage{lmodern}
\usepackage{microtype}
\usepackage{amsmath,amssymb,amsthm}
\usepackage{algorithm}
\usepackage{algpseudocode}
\usepackage{abstract}                 % gives us \begin{onecolabstract}
\usepackage{booktabs}
\usepackage{multirow}
\usepackage{graphicx}
\usepackage{xcolor}
\usepackage[hidelinks,colorlinks=true,linkcolor=blue!50!black,citecolor=blue!50!black,urlcolor=blue!50!black]{hyperref}
\usepackage{url}
\usepackage{titlesec}
\usepackage{caption}
\usepackage[numbers,sort&compress]{natbib}

% Tighter section spacing for a two-column layout.
\titlespacing*{\section}{0pt}{8pt}{4pt}
\titlespacing*{\subsection}{0pt}{6pt}{3pt}
\captionsetup{font=small,labelfont=bf}

% Custom math operators / shortcuts.
\newcommand{\INV}[1]{\textsc{Inv-#1}}
\newcommand{\Risk}{\mathrm{risk}}
\newcommand{\Judge}{\mathcal{J}}
\newcommand{\Cand}{n_{\text{cand}}}
\newcommand{\Reuse}{u}
\newcommand{\Shuf}{s}
\newcommand{\GateDecision}{\textsc{GateDecision}}
\newcommand{\code}[1]{\texttt{\small #1}}

% Bibliography style: numeric IEEE-like.
\bibliographystyle{IEEEtran}

% Theorem environment for Theorem 1 (Zero Violations).
\newtheorem{theorem}{Theorem}

% ----------------------------------------------------------------------------
% Title block
% ----------------------------------------------------------------------------
\title{\textbf{\textsc{Inv-15}}: A Formal Safety Invariant for KV-Cache
       Reuse in Multi-Agent Judge Pipelines}

\author{Pablo M. Suarez\thanks{Correspondence:
        \href{mailto:suarezpm@csnat.unt.edu.ar}{suarezpm@csnat.unt.edu.ar}
        (academic) \textperiodcentered{}
        \href{mailto:p.ms.08@hotmail.com}{p.ms.08@hotmail.com}.
        Code, benchmarks, and DevCloud-ATL1 logs:
        \href{https://github.com/SuarezPM/Apohara\_Context\_Forge}%
             {github.com/SuarezPM/Apohara\_Context\_Forge}.
        Author: \href{https://github.com/SuarezPM}{GitHub @SuarezPM}.
        \textbf{DOI}: \href{https://doi.org/10.5281/zenodo.20114594}%
             {10.5281/zenodo.20114594} (Zenodo, indexed by Google Scholar).}\\
        APOHARA \textperiodcentered{} ContextForge}

\date{May 10, 2026 \textperiodcentered{}
      \href{https://doi.org/10.5281/zenodo.20114594}%
           {DOI: 10.5281/zenodo.20114594}}

\begin{document}
\twocolumn[
  \maketitle
  \begin{onecolabstract}
    \noindent
    Cross-agent KV-cache reuse is the single largest optimization available to
    multi-agent large language model (LLM) pipelines\textemdash systems in
    which a chain of agents (retriever, reranker, summarizer, critic,
    responder) shares a common context. Recent work
    \cite{kvcomm,tokendance,lmcache} demonstrates compression ratios of
    7--17$\times$ when KV blocks are deduplicated across agents. However,
    reuse silently degrades \emph{judge} and \emph{critic} agents: when the
    judge compares multiple candidates, attention patterns cached from a
    prior ranking corrupt its independence. This failure mode\textemdash
    \emph{Judge Candidate Reuse} (JCR)\textemdash was identified
    theoretically in \cite{jcrfailure} but has not been resolved by any
    production KV-coordination system.

    We introduce \INV{15}, a formal safety invariant requiring that any
    judge-type agent whose JCR risk score exceeds a fixed threshold use
    \emph{dense} prefill\textemdash bypassing the shared KV registry. The risk
    score is a closed-form function of agent role, candidate count, reuse
    rate, and candidate-layout shuffling. We implement \INV{15} as the
    \emph{JCR Safety Gate} in ContextForge, an AMD-native KV-coordination
    layer running on Instinct MI300X (ROCm 7.x) through the vLLM V1 ATOM
    plugin. On an end-to-end 15-scenario benchmark with the
    Qwen3.6-235B-A22B mixture-of-experts model on AMD DevCloud ATL1, we
    observe \textbf{zero \INV{15} violations} across the full sweep, a
    critic dense-prefill rate of \textbf{$1.000$}, and full compatibility
    with $10.81\times$ TokenDance compression\textemdash with no measurable
    overhead beyond the gate's $5\,\mathrm{ms}$ decision latency. To our
    knowledge, this is the first production implementation of a formal
    safety invariant for cross-agent KV-cache reuse. We argue \INV{15}
    should be a standard preflight check in any system that shares KV state
    across judge-class agents.

    \vspace{0.5em}
    \noindent\textbf{Keywords:} multi-agent LLMs, KV-cache, system
    invariants, judge consistency, vLLM, AMD MI300X, ROCm.
    \vspace{1em}
  \end{onecolabstract}
]

% ============================================================================
\section{Introduction}
\label{sec:intro}
% ============================================================================

Multi-agent large language model (LLM) pipelines have become a standard
deployment pattern for retrieval-augmented generation, agentic reasoning, and
LLM-as-a-judge evaluation
\cite{kvcomm,tokendance,specdecode}. A
canonical pipeline\textemdash $\textsc{Retriever} \rightarrow
\textsc{Reranker} \rightarrow \textsc{Summarizer} \rightarrow \textsc{Critic}
\rightarrow \textsc{Responder}$\textemdash invokes a 35--235B parameter model
five times in sequence, each agent consuming a context that overlaps
substantially (often by $40$--$60\,\%$) with the others. On the AMD Instinct
MI300X accelerator's $192\,\mathrm{GB}$ HBM3, this redundancy materializes
the same KV-cache entries five times: the shared system prompt, the user
query, and the retrieved documents. Without coordination, more than half of
the GPU's memory budget is wasted before the first output token is
generated.

The optimization is well-studied. \emph{vLLM Automatic Prefix Caching} (APC)
deduplicates KV blocks within a single tenant
\cite{vllm}; LMCache extends caching across nodes \cite{lmcache};
KVCOMM \cite{kvcomm} introduces SimHash anchor matching for cross-context
offset hints; TokenDance \cite{tokendance} stores one master KV cache plus
$N-1$ sparse diffs, achieving $11$--$17\times$ compression in committee
inference. Each of these systems treats the cache as an unstructured pool of
attention state, freely reusable wherever a content match exists.

This assumption breaks for one specific agent role: the \emph{judge}. When
a critic agent evaluates a set of candidates $\{c_1, c_2, \ldots, c_k\}$,
the cached attention patterns from a prior evaluation encode the prior
candidate ordering. If the cache is reused\textemdash even partially\textemdash
the judge's verdict is no longer independent: it is conditioned on
positional priors from a different evaluation. The recent theoretical work
of \cite{jcrfailure} formalizes this failure mode as \emph{Judge Candidate
Reuse} (JCR) and shows empirically that judge consistency rates degrade by
$8$--$23\,\%$ under aggressive cross-agent KV reuse, without any other
metric (latency, throughput, accuracy on non-judge tasks) signalling the
regression. Crucially, \cite{jcrfailure} does \emph{not} propose a
production-grade mitigation; the authors leave it as future work.

\textbf{The gap.} Existing KV-coordination systems share three properties:
(1) they optimize purely for throughput or memory; (2) they have no notion
of agent \emph{role}; (3) they expose no formal contract under which the
cache is safe to reuse. As deployments scale from research prototypes to
production agentic platforms, this gap becomes a correctness liability.

\textbf{Our contribution.} We propose \INV{15}, a formal safety invariant
that prohibits KV-cache reuse for any judge-class agent whose JCR risk score
exceeds a fixed threshold $\tau = 0.7$. \INV{15} is enforced by the
\emph{JCR Safety Gate}, a preflight check integrated into the multi-agent
pipeline ahead of KV-block materialization. The gate is implemented as part
of ContextForge \cite{contextforge}, an AMD-native KV-coordination layer
intercepting at the vLLM V1 ATOM plugin interface and running on Instinct
MI300X. We make four contributions:
\begin{enumerate}
  \item A formal definition of \INV{15} together with a closed-form risk
        score function over four observable features of the upcoming agent
        invocation.
  \item A production-grade gate algorithm (\GateDecision) with $O(1)$
        per-call cost and an audit-grade decision log.
  \item Empirical verification on real MI300X hardware: \emph{zero}
        \INV{15} violations across the full high-risk sweep, critic dense
        prefill rate $1.000$.
  \item A composition result: \INV{15} co-exists with $10.81\times$
        TokenDance compression in the same pipeline, demonstrating that
        safety and compression are not in tension.
\end{enumerate}

The remainder of this paper is organized as follows.
Section~\ref{sec:background} reviews cross-agent KV sharing and the JCR
failure mode. Section~\ref{sec:invariant} introduces \INV{15} and its risk
function. Section~\ref{sec:impl} describes the ContextForge implementation.
Section~\ref{sec:eval} reports benchmark results on MI300X. We discuss
extensions and limitations in Section~\ref{sec:discussion} and conclude in
Section~\ref{sec:conclusion}.

% ============================================================================
\section{Background}
\label{sec:background}
% ============================================================================

\subsection{KV-Cache Sharing in Multi-Agent Systems}
\label{sec:bg:kvshare}

In a transformer decoder, the key and value tensors of past tokens (the
\emph{KV cache}) dominate inference memory: for a 35B parameter model with
context length $L$, the per-layer KV footprint is
$2 \cdot L \cdot d_{\mathrm{kv}} \cdot b$ bytes, where $d_{\mathrm{kv}}$ is
the per-head KV dimension and $b$ is the data-type width
\cite{vllm,rotatekv}. With $L = 8192$ and FP16, a single context occupies
roughly $11.7\,\mathrm{GB}$.

In a multi-agent pipeline, the agents share a substantial prefix\textemdash
typically the system prompt, the user query, and any retrieved documents
\cite{kvcomm}. Without coordination, each agent re-materializes the same
prefix tensors independently. vLLM APC deduplicates within a single
inference engine via hash-keyed PagedAttention blocks. LMCache extends this
across processes and nodes via a content-addressable store
\cite{lmcache}. KVCOMM contributes \emph{cross-context offset hints}: a
SimHash signature is computed per block, near-duplicate blocks are reused
with a per-token RoPE offset correction, eliminating drift at the attention
layer.

TokenDance \cite{tokendance} represents the current state of the art for
\emph{committee} inference (multiple agents reasoning over the same
context). The KV cache is decomposed into a \emph{master} (the largest or
most authoritative agent's cache) plus $N-1$ sparse diffs against the
master. Empirically, when the agents share a long prefix and diverge only
in the role-specific suffix, the diff is mostly zero and a threshold-gated
sparse representation achieves $11$--$17\times$ compression with
reconstruction error below $10^{-4}$.

\subsection{The JCR Failure Mode}
\label{sec:bg:jcr}

\cite{jcrfailure} identify a failure mode unique to \emph{judge} agents.
A judge evaluates a set of candidates $\{c_1, \ldots, c_k\}$ and emits a
ranking, score, or verdict. Modern judge prompts present candidates in a
serialized order, e.g.:
\begin{quote}
\small
\textit{``Compare the following responses and select the most accurate.\\
Response A: \ldots \\ Response B: \ldots \\ Response C: \ldots''}
\end{quote}
The attention patterns the judge develops while reading these candidates
depend on \emph{positional} cues: which candidate appears first, how their
embeddings interact, and how the judge's hidden state evolves across the
sequence. If the KV cache from a \emph{previous} judge invocation is reused
(e.g., because the system prompt and rubric are identical), the cached
attention state encodes the prior ordering. The reused state biases the
new evaluation toward the prior verdict, even when the candidate texts
themselves are different.

\cite{jcrfailure} quantify this as the \emph{Judge Consistency Rate} (JCR),
the fraction of judge verdicts that are stable under permutation of
candidate order. They report JCR drops of $8$--$23\,\%$ under aggressive
KV reuse, with the largest drops occurring when (i) the number of
candidates is large, (ii) the candidate order is shuffled between
evaluations, and (iii) the reuse rate (fraction of KV blocks served from
cache) exceeds $\sim 80\,\%$. Critically, the authors observe no other
visible regression: latency, throughput, and non-judge accuracy remain
unchanged. The failure is silent.

\subsection{Compression vs.\ Safety}
\label{sec:bg:tension}

There is a structural tension between TokenDance-style sharing and judge
independence. TokenDance achieves its highest compression ratios precisely
when the master and mirror agents share long prefixes\textemdash the
condition under which the JCR failure is most likely. A naive deployment
that turns on TokenDance for all five agents in a Retriever-to-Responder
pipeline will silently corrupt the Critic's verdict whenever the candidate
layout changes between calls. The contribution of this paper is to resolve
that tension by gating the Critic\textemdash and only the
Critic\textemdash through a formal invariant.

% ============================================================================
\section{The JCR Safety Gate and \INV{15}}
\label{sec:invariant}
% ============================================================================

\subsection{Formal Definition}
\label{sec:invariant:def}

Let $\mathcal{A}$ denote the set of agent roles in the pipeline and
$\Judge \subseteq \mathcal{A}$ the set of \emph{judge-class} roles. In our
reference pipeline, $\Judge = \{\textsc{critic}, \textsc{judge}\}$. Let an
upcoming agent invocation be described by the tuple
\begin{equation*}
  \mathbf{x} = (\rho,\, \Cand,\, \Reuse,\, \Shuf) \in \mathcal{A}
                  \times \mathbb{Z}_{\geq 0}
                  \times [0,1] \times \{0,1\},
\end{equation*}
where $\rho$ is the agent role, $\Cand$ is the number of candidates the
agent will compare, $\Reuse \in [0,1]$ is the fraction of KV blocks the
shared registry would serve from cache absent any safety constraint, and
$\Shuf \in \{0,1\}$ indicates whether the candidate layout has changed
since the agent's previous invocation. Define the \emph{JCR risk score}
$\Risk : \mathcal{A} \times \mathbb{Z}_{\geq 0} \times [0,1] \times \{0,1\}
\rightarrow [0,1]$ as:
\begin{align}
  \Risk(\mathbf{x}) \;=\; \min\Big(1,\;
    & b(\rho)
      + \alpha_n\,\mathbb{1}[\Cand \geq 3] \nonumber \\[-2pt]
    & + \alpha_s\,\Shuf
      + \alpha_u\,\mathbb{1}[\Reuse > 0.8]
    \Big),
  \label{eq:risk}
\end{align}
where
\begin{equation}
  b(\rho) =
  \begin{cases}
    b_J & \text{if } \rho \in \Judge \\
    b_O & \text{otherwise}
  \end{cases}
  \label{eq:base}
\end{equation}
with coefficients $b_J = 0.6$, $b_O = 0.1$, $\alpha_n = 0.2$,
$\alpha_s = 0.2$, $\alpha_u = 0.1$. The cap at $1.0$ ensures the score is
a proper risk index. We define \INV{15} as the following safety
contract:

\paragraph{\INV{15} (Judge Dense-Prefill Invariant).}
\emph{For every agent invocation $\mathbf{x}$ such that
$\rho \in \Judge$ and $\Risk(\mathbf{x}) > \tau$, KV-cache reuse is
prohibited; the invocation must use dense prefill.} We fix $\tau = 0.7$
(Section~\ref{sec:invariant:threshold}).

Equivalently, an invocation $\mathbf{x}$ \emph{satisfies} \INV{15} if and
only if
\begin{equation}
  \big(\rho \in \Judge \,\wedge\, \Risk(\mathbf{x}) > \tau\big)
  \;\Rightarrow\;
  \mathrm{prefill}(\mathbf{x}) = \textsc{dense}.
  \label{eq:inv15}
\end{equation}
A run \emph{violates} \INV{15} if any single invocation does not satisfy
the implication. A core property of our gate is that it produces zero
violations by construction (Theorem~\ref{thm:zero}).

\subsection{The Gate Algorithm}
\label{sec:invariant:algo}

Algorithm~\ref{alg:gate} specifies the \GateDecision\ procedure that
implements \INV{15}. The procedure runs once per agent invocation, before
the KV registry is queried.

\begin{algorithm*}[t]
\caption{\GateDecision: \INV{15}-enforcing preflight. A \textsc{true}
         return value signals that the invariant fires and the agent
         must use dense prefill. Cost is $O(1)$ per call.}
\label{alg:gate}
\begin{algorithmic}[1]
\Require agent role $\rho$, candidate count $\Cand$, reuse rate
         $\Reuse \in [0,1]$, layout-shuffled flag $\Shuf \in \{0,1\}$
\Require judge-role set $\Judge$, coefficients
         $(b_J, b_O, \alpha_n, \alpha_s, \alpha_u)$, threshold $\tau$
\Ensure \texttt{use\_dense\_prefill} $\in \{\textsc{true}, \textsc{false}\}$
\State $r \gets b_O$
\If{$\rho \in \Judge$} $r \gets b_J$ \EndIf
\If{$\Cand \geq 3$} $r \gets r + \alpha_n$ \EndIf
\If{$\Shuf = 1$} $r \gets r + \alpha_s$ \EndIf
\If{$\Reuse > 0.8$} $r \gets r + \alpha_u$ \EndIf
\State $r \gets \min(r, 1.0)$
\State $\texttt{log.append}\big((\rho, \Cand, \Reuse, \Shuf, r)\big)$
\If{$\rho \in \Judge$ \textbf{and} $r > \tau$}
  \State \Return \textsc{true}
\Else
  \State \Return \textsc{false}
\EndIf
\end{algorithmic}
\end{algorithm*}

\begin{theorem}[Zero violations]
\label{thm:zero}
Any pipeline that calls \GateDecision\ (Algorithm~\ref{alg:gate}) once per
agent invocation and respects its boolean output satisfies \INV{15} on
every invocation. Consequently the number of violations across any
sequence of invocations is exactly zero.
\end{theorem}

\begin{proof}[Proof sketch]
\GateDecision\ computes $\Risk(\mathbf{x})$ exactly as in
Eq.~\ref{eq:risk}, then returns \textsc{true} if and only if
$\rho \in \Judge$ and $r > \tau$. The pipeline's contract is to call
$\mathrm{prefill}(\mathbf{x}) = \textsc{dense}$ whenever the gate returns
\textsc{true}. Hence the implication of Eq.~\ref{eq:inv15} holds by
construction.
\end{proof}

\subsection{Integration at the ATOM Plugin Level}
\label{sec:invariant:atom}

The vLLM V1 architecture exposes the \texttt{vllm.general\_plugins}
entry-point (the \emph{ATOM plugin} surface) at which a third-party module
can intercept KV-block materialization \emph{before} attention is computed
\cite{vllm}. ContextForge registers a plugin that delegates every prefill
request through \GateDecision. When the gate returns \textsc{true}, the
plugin marks the request as \texttt{prefill\_mode=dense} on the
\texttt{SchedulerOutput}; the engine routes around the shared registry and
materializes a fresh KV cache for that agent's call only.

The cost of this interception is dominated by the dictionary lookup
$\rho \in \Judge$ plus a constant number of comparisons and additions; it
is $O(1)$ in the number of agents and is independent of context length.

\subsection{Choice of Threshold $\tau$}
\label{sec:invariant:threshold}

We fix $\tau = 0.7$ as the default. The choice can be justified on three
grounds:
\begin{enumerate}
  \item \textbf{Coverage.} Every combination of $(\rho \in \Judge,
        \Cand \geq 3, \Shuf = 1)$ produces $r \geq 0.6 + 0.2 + 0.2 = 1.0$,
        comfortably above $0.7$. Every combination of
        $(\rho \in \Judge, \Cand \geq 3)$ alone produces $r = 0.8 > 0.7$.
        Combinations with at most one risk factor (e.g.\ judge with two
        candidates and no shuffle) produce $r = 0.6 < 0.7$, falling
        through the gate.
  \item \textbf{Empirical alignment.} The high-risk regimes in
        \cite{jcrfailure} (3+ candidates, shuffled layout, or
        $>80\,\%$ reuse) correspond to JCR drops of $\geq 10\,\%$. The
        $\tau = 0.7$ contour selects exactly these regimes.
  \item \textbf{False-positive cost.} A false positive (firing dense
        prefill on a safe call) costs at most one re-materialization;
        a false negative (missing an unsafe call) silently corrupts the
        verdict. The asymmetry favors a conservative threshold.
\end{enumerate}
We treat $\tau$ as a deployment-time knob; Section~\ref{sec:discussion}
discusses adaptive thresholding.

% ============================================================================
\section{Implementation}
\label{sec:impl}
% ============================================================================

\subsection{ContextForge Architecture}
\label{sec:impl:arch}

ContextForge is an AMD-native KV-coordination layer that aggregates ten
peer-reviewed mechanisms\textemdash among them KVCOMM \cite{kvcomm},
TokenDance \cite{tokendance}, RotateKV \cite{rotatekv}, and the
queueing-aware vLLM scheduler of \cite{queueing}\textemdash behind a single
MCP server (FastAPI + asyncio). All KV operations flow through a
\texttt{ContextRegistry} that owns an LSH token matcher, a FAISS ANN index,
and a VRAM-aware cache with a five-mode pressure-responsive eviction
policy. \INV{15} sits at the entry of this pipeline as a preflight check;
no downstream component sees a request that violates the invariant.

\begin{figure}[t]
  \centering
  % If the figure asset is unavailable, the framed placeholder gives the
  % reader the structural picture.
  \fbox{\parbox{0.93\columnwidth}{
    \centering\small
    \textsc{Retriever} $\to$ \textsc{Reranker} $\to$
    \textsc{Summarizer} $\to$ \textsc{\textbf{Critic}}
    $\to$ \textsc{Responder} \\[2pt]
    \rule{0.85\columnwidth}{0.4pt} \\[2pt]
    All agents query \textbf{ContextRegistry} (LSH + FAISS + VRAM cache). \\
    \textbf{Critic} request first traverses \textsc{JCRSafetyGate}: \\
    if $\Risk > 0.7$, request is routed to \textbf{dense prefill}
    (bypasses registry).
  }}
  \caption{Placement of the JCR Safety Gate at the ATOM plugin entry. The
  gate sits ahead of TokenDance and the registry so that no shared block is
  ever returned to a judge invocation that fails \INV{15}.}
  \label{fig:arch}
\end{figure}

\subsection{The JCRSafetyGate Class}
\label{sec:impl:class}

The gate is exposed as the Python class
\code{apohara\_context\_forge.safety.jcr\_gate.JCRSafetyGate}. Public API:

\begin{itemize}
  \item \code{compute\_jcr\_risk(role, n, u, s) -> float}\textemdash
        returns $\Risk(\mathbf{x})$ in $[0, 1]$.
  \item \code{should\_use\_dense\_prefill(role, n, u, s) -> bool}\textemdash
        the boolean form of \INV{15}.
  \item \code{gate\_decision(role, n, u, s) -> JCRDecision}\textemdash
        records an audit log entry with role, risk, decision, and
        human-readable reason.
  \item \code{summary() -> dict}\textemdash aggregate decision counts,
        critic dense rate, mean risk over the run.
\end{itemize}

The audit log is the means by which a deployment proves \INV{15}
compliance after the fact: every gate decision is recorded with the four
input features, the computed risk, the boolean output, and a timestamp.
Compliance audits scan the log for the predicate
$\rho \in \Judge \wedge r > \tau \wedge \mathrm{use\_dense} = \textsc{false}$;
the count of matching entries is, by Theorem~\ref{thm:zero}, exactly zero.

\subsection{Composition with TokenDance}
\label{sec:impl:tokendance}

TokenDance \cite{tokendance} stores one master KV cache plus $N-1$ sparse
diffs. We register the gate \emph{before} the TokenDance reconstruction
step: a judge invocation that fails \INV{15} never reaches the master, so
the dense prefill is computed independently. Non-judge invocations
(retriever, reranker, summarizer, responder) flow through TokenDance as
usual. Empirically the composition retains the full $10.81\times$
compression ratio (Section~\ref{sec:eval:tokendance}).

\subsection{AMD MI300X Specifics}
\label{sec:impl:mi300x}

ContextForge targets the AMD Instinct MI300X (CDNA3, $192\,\mathrm{GB}$
HBM3, 304 compute units) via ROCm 7.x and the AITER kernel library
\cite{aiter}. The gate is pure Python and pure CPU; it does not consume any
GPU compute. Its only platform-specific dependency is the
\texttt{AITER\_ENABLE\_VSKIP=0} environment flag, which we hard-code via
\code{AITERConfig.apply()} to prevent a known kernel crash in the fused
MoE path triggered when dense-prefill traffic is interleaved with shared
prefill on the same scheduler step.

% ============================================================================
\section{Evaluation}
\label{sec:eval}
% ============================================================================

\subsection{Experimental Setup}
\label{sec:eval:setup}

We run all experiments on AMD DevCloud ATL1, on a single Instinct MI300X
($192\,\mathrm{GB}$ HBM3, ROCm 7.x). The inference engine is vLLM V1 with
the ContextForge ATOM plugin enabled. The model under test is
Qwen3.6-235B-A22B, a $235$B-parameter mixture-of-experts decoder with
$22$B active parameters per token; this is the largest open MoE that fits
end-to-end on a single MI300X with sufficient headroom for the
$5$-agent pipeline. The reference pipeline is
$\textsc{Retriever} \to \textsc{Reranker} \to \textsc{Summarizer} \to
\textsc{Critic} \to \textsc{Responder}$, with the Critic in
$\Judge = \{\textsc{critic}, \textsc{judge}\}$. All measurements are from
a single benchmark run completed on 2026-05-10; the full log is reproducible
from the repository's \code{logs/benchmark\_v6\_final.txt}.

\subsection{Benchmark Sweep}
\label{sec:eval:sweep}

We construct a $9$-point sweep over the risk axes: $5$ high-risk
combinations (Critic role, $\Cand \in \{3, 4, 5, 6\}$,
$\Shuf \in \{0, 1\}$, $\Reuse \in \{0.5, 0.85, 0.9, 0.95\}$) and $4$
low-risk combinations (non-judge roles at the same extremes). For each,
the pipeline issues a single agent invocation and we record (a) the gate
decision, (b) the risk score, and (c) whether the implication of
Eq.~\ref{eq:inv15} held. Results are reported in Table~\ref{tab:sweep}.

\begin{table}[t]
  \centering
  \small
  \caption{Benchmark sweep results (S-15
           \texttt{jcr\_gate\_critic\_safety}). Critic-class invocations
           with at least one risk factor produce $\Risk = 1.0$ (capped) and
           fire \INV{15}. Non-judge invocations never fire the gate even
           at extreme settings. Zero violations on $9/9$ decisions.}
  \label{tab:sweep}
  \begin{tabular}{@{}llrrrl@{}}
    \toprule
    Role & $\Cand$ & $\Shuf$ & $\Reuse$ & $\Risk$ & Dense? \\
    \midrule
    \multicolumn{6}{l}{\textit{High-risk (Critic)}} \\
    \quad critic & 5 & 1 & 0.90 & 1.000 & \textbf{yes} \\
    \quad critic & 4 & 1 & 0.85 & 1.000 & \textbf{yes} \\
    \quad critic & 3 & 1 & 0.95 & 1.000 & \textbf{yes} \\
    \quad critic & 5 & 1 & 0.50 & 1.000 & \textbf{yes} \\
    \quad critic & 6 & 0 & 0.85 & 0.900 & \textbf{yes} \\
    \midrule
    \multicolumn{6}{l}{\textit{Low-risk (non-judge)}} \\
    \quad retriever  & 2 & 1 & 0.90 & 0.400 & no \\
    \quad reranker   & 5 & 1 & 0.95 & 0.600 & no \\
    \quad summarizer & 3 & 0 & 0.90 & 0.400 & no \\
    \quad responder  & 5 & 1 & 0.80 & 0.500 & no \\
    \midrule
    \multicolumn{5}{r}{\INV{15} violations:} & \textbf{0 / 9} \\
    \bottomrule
  \end{tabular}
\end{table}

\subsection{\INV{15} Violation Rate}
\label{sec:eval:inv}

Across the $9$-point sweep, \INV{15} fires for every Critic invocation
with risk $> 0.7$ (all $5$ high-risk cases) and never fires for
non-judge invocations (all $4$ low-risk cases). The number of
violations\textemdash defined as a Critic invocation with $r > 0.7$ for
which \texttt{use\_dense\_prefill = false}\textemdash is zero. This is
the expected outcome of Theorem~\ref{thm:zero}; we report it as
empirical confirmation rather than as a discovery.

\subsection{Critic Dense Prefill Rate}
\label{sec:eval:critic}

We define the \emph{critic dense prefill rate} as the fraction of Critic
decisions for which the gate returned \textsc{true}. Over the $5$ Critic
invocations in the sweep this rate is $1.000$ (target $\geq 0.5$). All
five Critic invocations had at least one risk factor and triggered
\INV{15}.

\subsection{Gate Overhead}
\label{sec:eval:overhead}

The gate's per-decision cost is bounded by a small constant number of
dictionary lookups and arithmetic operations on Python floats. We measure
total time for the $9$-point sweep at $5\,\mathrm{ms}$ wall-clock,
corresponding to $\sim\!1{,}800$ decisions per second. This is three
orders of magnitude below the time required to materialize a single KV
block on the MI300X; the gate is therefore not a throughput bottleneck.

\subsection{Compatibility with TokenDance}
\label{sec:eval:tokendance}

To validate the safety/compression composition, we run S-14
(\code{token\_dance\_compression}) on a $12$-agent committee: one master
plus $11$ mirrors. The master holds a $200$-block KV cache; each mirror
diverges in at most two blocks. We measure the post-compression footprint
with \INV{15} enabled (the Critic agent runs through dense prefill, all
others share via TokenDance). The achieved compression ratio is
$\mathbf{10.81\times}$ (target $\geq 10\times$); reconstruction error is
$1.19 \times 10^{-7}$ (target $\leq 1 \times 10^{-4}$). Enabling
\INV{15} has no measurable effect on the TokenDance ratio because the
Critic's KV cache is a small fraction of the total committee footprint
and was already a high-divergence agent. See Table~\ref{tab:e2e}.

\subsection{End-to-End Pipeline}
\label{sec:eval:e2e}

The full $15$-scenario benchmark (V4 baseline + V5 systems + V6 safety)
passes $15/15$. The live $5$-agent demo on the same hardware reports
\textbf{$79.85\,\%$ token savings} ($263 \to 53$ tokens), an average
TTFT of $23.78\,\mathrm{ms}$ per agent, and the Critic correctly routed
through dense prefill. The unit test suite of ContextForge passes
$310/310$ with zero failures.

\begin{table}[t]
  \centering
  \small
  \caption{End-to-end metrics on MI300X (DevCloud ATL1, 2026-05-10).
           Numbers in bold are the \INV{15}-relevant subset.}
  \label{tab:e2e}
  \begin{tabular}{@{}lr@{}}
    \toprule
    Metric & Value \\
    \midrule
    Benchmark scenarios passed                  & $15 / 15$ \\
    Unit tests passed                           & $310 / 310$ \\
    Failed tests                                & $0$ \\
    \midrule
    \textbf{\INV{15} violations (sweep)}        & $\mathbf{0}$ \\
    \textbf{Critic dense prefill rate}          & $\mathbf{1.000}$ \\
    \textbf{Gate throughput}                    & $\mathbf{1{,}800\,\mathrm{decisions/s}}$ \\
    \textbf{Mean gate latency}                  & $\mathbf{0.56\,\mathrm{ms}}$ \\
    \midrule
    TokenDance compression                      & $10.81\times$ \\
    TokenDance reconstruction error             & $1.19\times10^{-7}$ \\
    Live token savings (5-agent demo)           & $79.85\,\%$ \\
    Mean TTFT (5-agent demo)                    & $23.78\,\mathrm{ms}$ \\
    \bottomrule
  \end{tabular}
\end{table}

% ============================================================================
\section{Discussion}
\label{sec:discussion}
% ============================================================================

\paragraph{When is \INV{15} too conservative?}
The gate is conservative by design: any false positive costs one
re-materialization, while any false negative silently corrupts a judge
verdict. Under workloads where the Critic is invoked frequently with the
same candidate set and the same layout (e.g., deterministic preference
benchmarks), \INV{15} will fire on every invocation and forfeit the
TokenDance gains for that one agent. Empirically this is a small price:
the Critic's KV footprint is a fraction of the pipeline's, and dense
prefill on a $\sim\!2\,\mathrm{k}$-token critic context completes in tens
of milliseconds on MI300X. We have not observed a workload in which this
overhead is the binding constraint, but we acknowledge it as a tunable.

\paragraph{Extension to other agent roles.}
\INV{15} as stated covers $\{\textsc{critic}, \textsc{judge}\}$. The same
risk model extends naturally to \emph{verifier} and \emph{planner} agents
whose outputs are sensitive to candidate ordering or world-state caching.
For verifiers, the analog of $\Shuf$ is whether the verified statement
has changed since the last invocation; for planners, the analog is whether
the goal stack has been re-shuffled. Generalizing the role set requires
re-tuning the coefficients $(b_J, \alpha_n, \alpha_s, \alpha_u)$, but the
algorithmic structure of Algorithm~\ref{alg:gate} is unchanged.

\paragraph{Adaptive thresholding.}
A natural follow-up is to learn $\tau$ from observed JCR rates: if the
deployment exposes a ground-truth JCR signal (e.g., from human raters or
held-out reference verdicts), $\tau$ can be tuned to minimize the
joint loss
$\lambda_{\mathrm{FP}}\cdot \mathrm{FP}(\tau) +
 \lambda_{\mathrm{FN}}\cdot \mathrm{FN}(\tau)$
where the asymmetry of the loss weights $\lambda_{\mathrm{FP}} \ll
\lambda_{\mathrm{FN}}$ reflects the asymmetric cost. We leave this to
future work.

\paragraph{Limitations.}
Our risk score is a closed-form heuristic, not a learned model. It does
not capture interactions among the four features (e.g., shuffled layout
may matter less for $\Cand = 3$ than for $\Cand = 6$). It also does not
distinguish between \emph{candidate} reuse (the failure mode of
\cite{jcrfailure}) and \emph{system-prompt} reuse, which is generally
safe. A more refined gate would separate these dimensions; we leave that
to future work. Finally, our benchmark is single-node single-GPU; the
multi-node extension via LMCache \cite{lmcache} will require additional
gate-state replication for which the audit log already provides the
groundwork.

% ============================================================================
\section{Conclusion}
\label{sec:conclusion}
% ============================================================================

We presented \INV{15}, a formal safety invariant for KV-cache reuse in
multi-agent LLM pipelines, and the JCR Safety Gate that enforces it. The
invariant is closed-form, deterministically computable in $O(1)$, and
provably violation-free by construction. We implemented the gate as part
of ContextForge, ran it on AMD Instinct MI300X with vLLM V1, and observed
zero violations and a $1.000$ critic dense prefill rate across a full
high-risk sweep, while preserving $10.81\times$ TokenDance compression
for the non-judge agents. To our knowledge this is the first production
implementation of a formal safety contract for cross-agent KV-cache reuse.
We argue that any future multi-agent KV-coordination
layer\textemdash whether targeted at TokenDance-style sharing, KVCOMM
anchor matching, or LMCache distributed caching\textemdash should include
a preflight gate of this form. Safe KV reuse is not at odds with
aggressive compression; it is a prerequisite for it.

\paragraph{Reproducibility.}
The complete implementation, benchmark scenarios, and DevCloud-ATL1 logs
are released under Apache-2.0 at
\href{https://github.com/SuarezPM/Apohara_Context_Forge}%
     {github.com/SuarezPM/Apohara\_Context\_Forge}; the safety gate is at
\code{apohara\_context\_forge/safety/jcr\_gate.py} and the benchmark
scenario is S-15 in \code{demo/benchmark\_v5.py}.

% ----------------------------------------------------------------------------
\bibliography{references}
% ----------------------------------------------------------------------------

\end{document}
